% Wersja źródłowa: audyt II, jednostka 26
\setcounter{section}{25}
\section{Geometria przestrzeni Banacha: refleksywność, zwartość słaba i wypukłość jednostajna}
\label{sec:geometria26}

\providecommand{\atlasref}[2]{#2}


\noindent\textbf{Wyniki wykorzystywane.} Korzystamy z twierdzenia
Hahna--Banacha dla przestrzeni unormowanych
(Twierdzenie~\ref{thm:hb-normed}), separacji wypukłej
(Wniosek~\ref{cor:point-closed-convex}), podstaw topologii
słabych z rozdziału~8, twierdzenia Banacha--Alaoglu
(Wniosek~\ref{cor:alaoglu-norm}) oraz zasady jednostajnej
ograniczoności (Wniosek~\ref{cor:uniform-boundedness}).
Twierdzenia Goldstine'a, Eberleina--\v Smuliana i Milmana--Pettisa
udowodnimy tutaj.
\par\medskip

W tym rozdziale $X$ oznacza przestrzeń Banacha nad $\K\in\{\R,\C\}$,
$X'$ jej przestrzeń dualną topologiczną z normą operatorową,
a $X''=(X')'$ i $X'''=(X'')'$ kolejne przestrzenie dualne
z tymi samymi konwencjami. Funkcjonały są liniowe nad $\K$,
również w przypadku zespolonym.
Przez $B_X=\{x\in X:\|x\|\le1\}$ rozumiemy domkniętą kulę
jednostkową, natomiast $\sigma(X,X')$ oznacza topologię słabą.
Domknięcia z indeksem $w$ dotyczą topologii słabej,
a z indeksem $w^\ast$ w $X''$ topologii $\sigma(X'',X')$.
Celem jest ustalenie, kiedy ograniczoność daje zwartość słabą,
co wolno wówczas wnioskować o ciągach oraz jak geometria normy
wpływa na refleksywność.

\subsection{Kanoniczne zanurzenie i twierdzenie Goldstine'a}

\noindent\emph{Powiązania tematyczne w Bibliotece matematycznej:} \atlasref{LA-07::thm:lin7-bidualnosc}{LA-07: kanoniczne odwzorowanie w podwójną przestrzeń sprzężoną w wymiarze skończonym}.


\begin{definicja}[Refleksywność]
\label{def:refleksywnosc26}
Przestrzeń Banacha $X$ nazywamy \emph{refleksywną}, jeżeli
kanoniczne odwzorowanie liniowe
$J_X:X\to X''$, określone wzorem
$(J_Xx)(f)=f(x)$ dla $x\in X$ i $f\in X'$, jest surjektywne.
\end{definicja}

\begin{stwierdzenie}[Kanoniczne zanurzenie]
\label{prop:canonical26}
Dla każdej przestrzeni unormowanej $X$ odwzorowanie
$J_X:X\to X''$ jest liniową izometrią. Jeżeli $X$ jest zupełna,
to $J_X(X)$ jest normowo domkniętą podprzestrzenią $X''$.
Ponadto $J_X$ jest homeomorfizmem przestrzeni
$(X,\sigma(X,X'))$ na swój obraz z topologią podprzestrzeni
$(X'',\sigma(X'',X'))$.
\end{stwierdzenie}
\begin{proof}
Dla ustalonego $x\in X$ odwzorowanie $f\mapsto f(x)$
jest liniowe oraz $|f(x)|\le\|f\|\|x\|$, więc należy
do $X''$. Liniowość $J_X$ wynika z liniowości każdego $f$.
Dla $x\in X$ mamy
$\|J_Xx\|=\sup_{\|f\|\le1}|f(x)|\le\|x\|$.
Dla $x\ne0$ Wniosek~\ref{cor:norming} zapewnia funkcjonał
$f\in X'$ o normie $1$ i wartości $f(x)=\|x\|$.
Daje to odwrotną nierówność; dla $x=0$ mamy równość.
Izometria jest różnowartościowa. Obraz przestrzeni zupełnej
przez izometrię jest zupełny, a zupełna podprzestrzeń
przestrzeni metrycznej jest domknięta.

Topologia słaba na $X$ jest wyznaczona przez seminormy
$x\mapsto|f(x)|$ dla $f\in X'$, natomiast topologia
słaba-* na $X''$ przez $\Phi\mapsto|\Phi(f)|$ dla
tych samych $f$. Ponieważ $|(J_Xx)(f)|=|f(x)|$,
przeciwobrazy tych seminorm są dokładnie rodziną
wyznaczającą topologię słabą. Stąd $J_X$ i jego
odwrotność na obrazie są ciągłe w podanych topologiach.
\end{proof}

\begin{twierdzenie}[Goldstine]
\label{thm:goldstine26}
Niech $X$ będzie przestrzenią unormowaną nad $\K$.
Wówczas $J_X(B_X)$ jest gęsty w $B_{X''}$ względem topologii
$\sigma(X'',X')$.
\end{twierdzenie}
\begin{proof}
Z izometryczności $J_X$ mamy $J_X(B_X)\subseteq B_{X''}$.
Kula $B_{X''}$ jest słabo-* domknięta, ponieważ jest
przecięciem zbiorów $\{\Psi:|\Psi(f)|\le1\}$ po
$f\in B_{X'}$. Przypuśćmy, że
$\Phi\in B_{X''}\setminus\overline{J_X(B_X)}^{\,w^\ast}$.
Z bazy otoczeń słabych-* (Twierdzenie~\ref{thm:slaba-baza})
istnieją $m\ge1$, $f_1,\ldots,f_m\in X'$ i $\eta>0$,
dla których żaden $x\in B_X$ nie spełnia wszystkich
nierówności $|f_j(x)-\Phi(f_j)|<\eta$.
Wektor $(\Phi(f_1),\ldots,\Phi(f_m))$ leży więc poza
domknięciem w $\K^m$ wypukłego zbioru
$\{(f_1(x),\ldots,f_m(x)):x\in B_X\}$.
Stosując Wniosek~\ref{cor:point-closed-convex} w
skończonym wymiarze, oddzielamy ten wektor od domknięcia.
Nad $\C$ rzeczywisty funkcjonał liniowy na $\C^m$
jest częścią rzeczywistą funkcjonału zespolonego
(Lemat~\ref{lem:realification}); nad $\R$ przyjmujemy
$\operatorname{Re}t=t$. Otrzymujemy skalary
$c_1,\ldots,c_m\in\K$ i liczbę rzeczywistą $a$ spełniające
\[
\sup_{x\in B_X}\operatorname{Re}\Big(\sum_{j=1}^m c_jf_j(x)\Big)
<a<\operatorname{Re}\Big(\sum_{j=1}^m c_j\Phi(f_j)\Big).
\]
Dla $f=\sum_jc_jf_j$ mamy
$\sup_{x\in B_X}\operatorname{Re}f(x)=\|f\|$:
nierówność $\le$ wynika z definicji normy, a dla
każdego $x$ z $f(x)\ne0$ można pomnożyć $x$ przez
skalar modułu $1$, aby otrzymać wartość $|f(x)|$.
W przypadku $f(x)=0$ nie trzeba zmieniać $x$.
Z liniowości $\Phi$ prawa strona nierówności to
$\operatorname{Re}\Phi(f)\le|\Phi(f)|\le\|\Phi\|\|f\|
\le\|f\|$. Otrzymana sprzeczność kończy dowód.
\end{proof}

\begin{twierdzenie}[Refleksywność a zwartość słaba kuli]
\label{thm:reflexive-ball26}
Niech $X$ będzie przestrzenią Banacha. Następujące warunki
są równoważne:
\begin{enumerate}[label=\textup{(\roman*)}]
\item $X$ jest refleksywna;
\item $B_X$ jest zwarta w topologii $\sigma(X,X')$;
\item każdy niepusty, ograniczony, normowo domknięty i wypukły
podzbiór $X$ jest słabo zwarty.
\end{enumerate}
\end{twierdzenie}
\begin{proof}
Jeżeli $J_X$ jest surjektywne, to z izometryczności i
Stwierdzenia~\ref{prop:canonical26} jego ograniczenie
do $B_X$ jest homeomorfizmem na $B_{X''}$ w odpowiednich
topologiach słabej i słabej-*. Stosujemy Wniosek~
\ref{cor:alaoglu-norm} do przestrzeni unormowanej $X'$:
kula jej przestrzeni dualnej $X''$ jest słabo-* zwarta.
Zatem (i) pociąga (ii).
Jeżeli zachodzi (ii), to dla każdego ograniczonego normowo,
zamkniętego i wypukłego $C\subseteq X$ istnieje $R>0$ z
$C\subseteq RB_X$. Z Wniosku~\ref{cor:weak-convex-closure} wynika, że
normowo domknięty zbiór wypukły jest słabo domknięty.
Zatem $C$ jest domkniętym podzbiorem słabo zwartej kuli $RB_X$,
co daje (iii). Przypadek $C=\varnothing$ nie wpływa na argument.
Warunek (iii) zastosowany do $B_X$ daje (ii).
Wreszcie jeśli $B_X$ jest słabo zwarta, to $J_X(B_X)$ jest
słabo-* zwarta, a zatem słabo-* domknięta w $X''$,
gdyż funkcjonały ewaluacji na $X'$ rozdzielają punkty
$X''$. Twierdzenie~\ref{thm:goldstine26} daje
$J_X(B_X)=B_{X''}$. Dla $\Phi\ne0$ wybieramy
$x\in B_X$ takie, że $J_Xx=\Phi/\|\Phi\|$; wtedy
$J_X(\|\Phi\|x)=\Phi$. Element zerowy jest obrazem zera.
Otrzymujemy surjektywność $J_X$, czyli (i).
\end{proof}

\begin{wniosek}[Podprzestrzenie i przestrzeń dualna]
\label{cor:reflexivity-stability26}
Jeżeli $X$ jest refleksywną przestrzenią Banacha, to każda
jej domknięta podprzestrzeń liniowa jest refleksywna.
Ponadto $X$ jest refleksywna wtedy i tylko wtedy, gdy $X'$
jest refleksywna.
\end{wniosek}
\begin{proof}
Niech $Y\subseteq X$ będzie domkniętą podprzestrzenią.
Jako zbiór domknięty i wypukły $Y$ jest słabo domknięta.
Zatem $B_Y=B_X\cap Y$ jest słabo zwarta w $X$.
Ograniczenia funkcjonałów z $X'$ należą do $Y'$,
a z Twierdzenia~\ref{thm:hb-normed} każdy $g\in Y'$
przedłuża się do $X'$ z zachowaniem normy.
Obie rodziny seminorm na $Y$ są więc identyczne,
co daje równość topologii słabej $Y$ i topologii
indukowanej z $X$. Ponieważ $Y$ jest normowo domknięta
w przestrzeni Banacha, sama jest przestrzenią Banacha. Twierdzenie~\ref{thm:reflexive-ball26}
daje refleksywność $Y$.

Jeżeli $X$ jest refleksywna, to $J_X:X\to X''$ jest
izomorfizmem izometrycznym. Dla $G\in X'''$ określamy
$f(x)=G(J_Xx)$. Liniowość oraz oszacowanie
$|f(x)|\le\|G\|\|x\|$ dają $f\in X'$.
Dla $\Phi=J_Xx\in X''$
mamy $(J_{X'}f)(\Phi)=\Phi(f)=f(x)=G(\Phi)$.
Surjektywność $J_X$ daje $J_{X'}f=G$, więc $X'$ jest refleksywna.
Odwrotnie, gdy $X'$ jest refleksywna i $J_X(X)\ne X''$,
obraz $J_X(X)$ jest normowo domknięty z zupełności $X$.
Twierdzenie~\ref{thm:point-closed-subspace}, zastosowane
do tej właściwej podprzestrzeni przestrzeni unormowanej $X''$,
daje niezerowy $G\in X'''$ zerujący się na $J_X(X)$.
Z refleksywności $X'$ wynika $G=J_{X'}f$ dla pewnego $f\in X'$.
Wtedy $0=G(J_Xx)=f(x)$ dla wszystkich $x\in X$, więc
$f=0$ i $G=0$, sprzeczność.
\end{proof}

\subsection{Twierdzenie Eberleina--\v Smuliana}

W dowolnej przestrzeni topologicznej zwartości nie można
charakteryzować wyłącznie ciągami. Dla słabej topologii
przestrzeni Banacha zachodzi jednak ważny wyjątek.

\begin{lemat}[Skończona rodzina normująca]
\label{lem:finite-norming26}
Niech $X$ będzie przestrzenią unormowaną,
a $M\subseteq X''$ skończenie wymiarową podprzestrzenią.
Istnieje skończony niepusty zbiór $F\subseteq B_{X'}$ taki, że
$\|\Psi\|\le 2\max_{f\in F}|\Psi(f)|$ dla wszystkich $\Psi\in M$.
\end{lemat}
\begin{proof}
Dla $M=\{0\}$ wystarczy przyjąć $F=\{0\}$.
W pozostałym przypadku dla każdego $\Psi\in M$ o normie $1$
wybieramy $f_\Psi\in B_{X'}$ z
$|\Psi(f_\Psi)|>3/4$, co jest możliwe z definicji
normy operatorowej $\Psi$. Sfera jednostkowa
skończenie wymiarowej przestrzeni $M$ jest zwarta
w topologii normowej.
Pokrywamy ją kulami $\{\Theta\in M:\|\Theta-\Psi\|<1/4\}$
i wybieramy skończone podpokrycie o środkach
$\Psi_1,\ldots,\Psi_r$. Dla $\Theta$ na sferze i odpowiedniego
$j$ mamy
\[
 |\Theta(f_{\Psi_j})|
 \ge|\Psi_j(f_{\Psi_j})|-\|\Theta-\Psi_j\|\|f_{\Psi_j}\|
 >3/4-1/4=1/2.
\]
Kładziemy $F=\{f_{\Psi_1},\ldots,f_{\Psi_r}\}$.
Dla niezerowego $\Theta\in M$ stosujemy nierówność
do $\Theta/\|\Theta\|$, a dla zera teza jest oczywista.
\end{proof}

\begin{twierdzenie}[Eberlein--\v Smulian]
\label{thm:eberlein-smulian26}
Niech $X$ będzie przestrzenią Banacha, a $A\subseteq X$
dowolnym podzbiorem. Następujące warunki są równoważne:
\begin{enumerate}[label=\textup{(\roman*)}]
\item domknięcie $\overline A^{\,w}$ jest słabo zwarte;
\item z każdego ciągu punktów $A$ można wybrać podciąg
zbieżny słabo do pewnego punktu $X$;
\item każdy ciąg w $A$ ma słaby punkt skupienia w $X$
(w każdym jego słabym otoczeniu leżą wyrazy ciągu o dowolnie dużych indeksach).
\end{enumerate}
\end{twierdzenie}
\begin{proof}
Gdy $A=\varnothing$, wszystkie warunki zachodzą.
Zakładamy więc $A\ne\varnothing$; przypadek $X=\{0\}$
jest natychmiastowy.

\emph{Krok pierwszy: (i) $\Rightarrow$ (ii).}
Niech $(x_n)\subseteq A$ i
$Y=\overline{\lin\{x_n:n\ge1\}}^{\,\|\cdot\|}$.
Jest to ośrodkowa podprzestrzeń: skończone kombinacje
liniowe $x_n$ o współczynnikach wymiernych (nad $\C$
o wymiernych częściach rzeczywistych i urojonych)
tworzą przeliczalny zbiór gęsty w $Y$.
Podprzestrzeń $Y$ jest normowo domknięta i wypukła,
więc jest słabo domknięta
(Wniosek~\ref{cor:weak-convex-closure}). Zbiór
$K=\overline A^{\,w}\cap Y$ jest słabo zwarty i zawiera $(x_n)$.
Wybieramy ciąg gęsty $(y_j)$ w $Y$ i funkcjonały
$f_j\in B_{X'}$ takie, że $f_j(y_j)=\|y_j\|$
(Wniosek~\ref{cor:norming}; dla $y_j=0$ bierzemy $f_j=0$).
Taka rodzina rozdziela punkty $Y$: dla $y\ne0$ wystarczy
wybrać $y_j$ w odległości mniejszej niż $\|y\|/3$ od $y$.
Wówczas $|f_j(y)|\ge\|y_j\|-\|y-y_j\|
\ge\|y\|-2\|y-y_j\|>\|y\|/3$.
Na $K$ funkcja
\[
d(u,v)=\sum_{j=1}^\infty 2^{-j}\min\{1,|f_j(u-v)|\}
\]
jest metryką: szereg jest zbieżny, symetria wynika
z modułu, a rozdzielanie punktów z powyższej własności
rodziny $(f_j)$. Nierówność trójkąta wynika z
podaddytywności $t\mapsto\min\{1,t\}$ dla $t\ge0$
oraz nierówności trójkąta dla modułu.
Dla ustalonego $v\in K$ każda suma częściowa
$d_N(\cdot,v)$ jest słabo ciągła i
$0\le d(u,v)-d_N(u,v)\le2^{-N}$ dla wszystkich $u,v\in K$.
Stąd $d(\cdot,v)$ jest słabo ciągła jako jednostajna
granica funkcji ciągłych; kule metryczne są słabo
otwarte w $K$. Topologia metryczna jest więc słabsza
od topologii słabej i jest Hausdorffa. Ponieważ $K$ jest słabo zwarty,
identyczność z $K$ słabo zwartą do $K$ metryczną jest
homeomorfizmem: ciągła bijekcja z przestrzeni zwartej
do Hausdorffa ma ciągłą odwrotność. W konsekwencji $K$
jest zwartą przestrzenią metryczną, więc $(x_n)$ ma w $K$
podciąg zbieżny słabo.

\emph{Krok drugi: (ii) $\Rightarrow$ (iii).}
Granica podciągu słabo zbieżnego jest słabym punktem skupienia
całego ciągu, gdyż indeksy podciągu rosną do nieskończoności.

\emph{Krok trzeci: (iii) $\Rightarrow$ (i).}
Najpierw $A$ jest ograniczony w normie. W przeciwnym razie
rodzina $\{J_Xx:x\in A\}\subseteq X''$, rozpatrywana jako
rodzina funkcjonałów na przestrzeni Banacha $X'$, byłaby
nieograniczona jednostajnie w normie, gdyż
$\|J_Xx\|=\|x\|$. Przestrzeń $X'$ jest przestrzenią Banacha nawet
dla unormowanej przestrzeni $X$. Z zasady jednostajnej
ograniczoności (Wniosek~\ref{cor:uniform-boundedness})
istniałby $f\in X'$ z
$\sup_{x\in A}|f(x)|=\infty$. Wybierając $x_n\in A$ z
$|f(x_n)|\ge n$, dostalibyśmy ciąg bez słabego punktu
skupienia: dla każdego $x\in X$ słabe otoczenie
$\{y:|f(y)-f(x)|<1\}$ nie zawiera żadnego $x_n$
dla $n>|f(x)|+1$.

Niech $\Phi\in\overline{J_X(A)}^{\,w^\ast}\subseteq X''$.
Skonstruujemy ciąg $(x_n)\subseteq A$ oraz skończone
zbiory $F_n\subseteq B_{X'}$. Wybieramy $x_1\in A$ i zbiór
$F_1$ z Lematu~\ref{lem:finite-norming26}, normujący z dokładnością
do czynnika $2$ przestrzeń
$M_1=\lin\{\Phi,J_Xx_1\}\subseteq X''$.
Gdy $x_1,\ldots,x_n$ i $F_1,\ldots,F_n$ są ustalone,
wybieramy $x_{n+1}\in A$ tak, by
$|\Phi(f)-f(x_{n+1})|<1/(n+1)$ dla każdego
$f\in F_1\cup\cdots\cup F_n$; jest to możliwe, ponieważ
$\Phi$ należy do słabo-* domknięcia zbioru $J_X(A)$,
a skończona rodzina warunków wyznacza
podstawowe otoczenie słabe-* punktu $\Phi$.
Następnie wybieramy $F_{n+1}$ normujący z dokładnością
do czynnika $2$ podprzestrzeń
$M_{n+1}=\lin\{\Phi,J_Xx_1,\ldots,J_Xx_{n+1}\}$.

Na mocy (iii) istnieje słaby punkt skupienia $x\in X$
ciągu $(x_n)$. Dla każdego $f\in\bigcup_nF_n$ mamy
$f(x_n)\to\Phi(f)$: jeśli $f\in F_j$, to
$|f(x_n)-\Phi(f)|<1/n$ dla $n>j$.
Punkt $f(x)$ jest punktem skupienia ciągu $(f(x_n))$,
bo przeciwobrazy otoczeń przez $f$ są słabo otwarte.
Jednoznaczność granicy skalarnej daje $f(x)=\Phi(f)$.
Ponadto $x$ należy do słabego domknięcia zbioru
$\{x_n\}$, a stąd do normowo domkniętej podprzestrzeni
$Y=\overline{\lin\{x_n:n\ge1\}}^{\,\|\cdot\|}$,
która jest słabo domknięta na mocy twierdzenia Hahna--Banacha.
W szczególności $J_Xx$ należy do domknięcia normowego
przestrzeni $\lin\{J_Xx_n:n\ge1\}$, zawartej w
$\bigcup_nM_n$. Przestrzenie $M_n$ są rosnące,
więc ich suma mnogościowa jest podprzestrzenią liniową.
Ponieważ $\Phi\in M_1$, również
$\Psi=\Phi-J_Xx$ należy do normowego domknięcia
$\bigcup_nM_n$. Dla dowolnego $\eta>0$ wybieramy
$n$ i $G\in M_n$ z $\|G-\Psi\|<\eta$. Z konstrukcji dla wszystkich
$f\in F_n$ zachodzi $\Psi(f)=0$, a więc
$G(f)=(G-\Psi)(f)$.
Ponieważ $G\in M_n$ oraz $F_n\subseteq B_{X'}$,
z własności rodziny normującej mamy
$\|G\|\le2\max_{f\in F_n}|G(f)|\le2\|G-\Psi\|$.
Stąd $\|\Psi\|\le3\|\Psi-G\|<3\eta$.
Dowolność $\eta>0$ daje $\Psi=0$, a zatem $\Phi=J_Xx$.

Pokazaliśmy, że słabo-* domknięcie $J_X(A)$ zawiera się
w $J_X(X)$. Z pierwszej części kroku trzeciego zbiór $A$
jest normowo ograniczony, więc istnieje $R>0$ takie,
że $J_X(A)\subseteq RB_{X''}$. Ta kula jest słabo-*
domknięta i zwarta (Wniosek~\ref{cor:alaoglu-norm}).
Zatem $C=\overline{J_X(A)}^{\,w^\ast}$ jest słabo-*
zwarty i, jak wykazaliśmy, $C\subseteq J_X(X)$.
Z własności domknięcia w podprzestrzeni wynika
$\overline{J_X(A)}^{\,w^\ast|_{J_X(X)}}=C\cap J_X(X)=C$.
Ponieważ $J_X$ jest homeomorfizmem na obraz
(Stwierdzenie~\ref{prop:canonical26}), otrzymujemy
$J_X(\overline A^{\,w})=C$.
Przenosząc zwartość przez odwrotność tego
homeomorfizmu, dostajemy słabą zwartość
$\overline A^{\,w}$.
\end{proof}

\begin{wniosek}[Ograniczone ciągi w przestrzeni refleksywnej]
\label{cor:bounded-subsequence26}
Jeżeli $X$ jest refleksywną przestrzenią Banacha,
to każdy ograniczony ciąg w $X$ ma podciąg słabo zbieżny.
\end{wniosek}
\begin{proof}
Wszystkie wyrazy ciągu należą do pewnej słabo zwartej kuli
$RB_X$ na mocy Twierdzenia~\ref{thm:reflexive-ball26}.
Stosujemy Twierdzenie~\ref{thm:eberlein-smulian26}.
\end{proof}

\subsection{Własność Schura i geometria przestrzeni \texorpdfstring{$\ell^1$}{ell1}}

\begin{definicja}[Własność Schura]
\label{def:schur26}
Przestrzeń Banacha $X$ ma \emph{własność Schura}, jeżeli dla
każdego ciągu $(x_n)\subseteq X$ i $x\in X$ ze zbieżności
$x_n\to x$ słabo wynika $\|x_n-x\|\to0$.
\end{definicja}

\begin{twierdzenie}[Własność Schura w $\ell^1$]
\label{thm:schur-l1-26}
Przestrzeń Banacha $\ell^1(\N,\K)$ ma własność Schura.
\end{twierdzenie}
\begin{proof}
Jeżeli $x_n\to x$ słabo, to $x_n-x\to0$ słabo,
więc wystarczy rozważyć ciąg $(x_n)$ słabo zbieżny
do zera. Dla każdej współrzędnej $j$ odwzorowanie
$x\mapsto x(j)$ jest liniowym funkcjonałem ciągłym,
gdyż $|x(j)|\le\|x\|_1$. Wobec tego $x_n(j)\to0$.
Przypuśćmy, że normy nie dążą do zera. Po przejściu do
podciągu możemy założyć, że $\|x_n\|_1\ge\varepsilon>0$
dla wszystkich $n$. Rekurencyjnie wybieramy podciąg
$y_k=x_{n_k}$ oraz liczby całkowite
$0=m_0<m_1<m_2<\cdots$ takie, że
\[
\sum_{j\le m_{k-1}}|y_k(j)|<\varepsilon/8,
\qquad
\sum_{j>m_k}|y_k(j)|<\varepsilon/8.
\]
Rozpoczynamy od $m_0=n_0=0$. Mając $m_{k-1}$ i
$n_{k-1}$, wybieramy $n_k>n_{k-1}$ tak duże, aby
$\sum_{j\le m_{k-1}}|x_{n_k}(j)|<\varepsilon/8$:
jest to skończona suma wyrazów dążących do zera
(dla $m_{k-1}=0$ suma jest zerowa).
Ustaliwszy $y_k=x_{n_k}$, wybieramy
$m_k>m_{k-1}$, korzystając z
$\sum_{j\ge1}|y_k(j)|<\infty$.
Na bloku $I_k=\{m_{k-1}+1,\ldots,m_k\}$ otrzymujemy
\[
 \sum_{j\in I_k}|y_k(j)|
 =\|y_k\|_1-\sum_{j\notin I_k}|y_k(j)|
 >3\varepsilon/4,\qquad
 \sum_{j\notin I_k}|y_k(j)|<\varepsilon/4.
\]
Bloki są parami rozłączne, a ponieważ $m_k\ge k$,
ich suma mnogościowa jest równa $\N$.
Definiujemy $a=(a_j)\in\ell^\infty$ tak, że dla $j\in I_k$
wybieramy $a_j=\overline{y_k(j)}/|y_k(j)|$, gdy $y_k(j)\ne0$,
i $a_j=0$ w przeciwnym wypadku.
Nad $\R$ sprzężenie jest identycznością.
Mamy $|a_j|\le1$ dla wszystkich $j$, a na $I_k$
zachodzi $a_jy_k(j)=|y_k(j)|$, także gdy $y_k(j)=0$.
Definiujemy jeden funkcjonał (niezależny od $k$)
$f(x)=\sum_{j\ge1}a_jx(j)$. Szereg jest bezwzględnie
zbieżny, gdyż $\sum_j|a_jx(j)|\le\|x\|_1$.
Liniowość wynika z liniowości sum częściowych,
a oszacowanie $|f(x)|\le\|x\|_1$ daje ciągłość,
zgodnie z Twierdzeniem~\ref{thm:dual-sequences12}.
Dla każdego $k$ z rozbicia bezwzględnie zbieżnego
szeregu i nierówności $\operatorname{Re}z\ge-|z|$
wynika
\[
\operatorname{Re}f(y_k)\ge
\sum_{j\in I_k}|y_k(j)|-\sum_{j\notin I_k}|y_k(j)|
>\varepsilon/2.
\]
Tymczasem $(y_k)$ jest podciągiem ciągu słabo
zbieżnego do zera, więc $f(y_k)\to0$ dla tego
ustalonego $f\in(\ell^1)'$. Otrzymujemy sprzeczność.
\end{proof}

\begin{wniosek}[Nierefleksywność przestrzeni $\ell^1$]
\label{cor:l1-nonreflexive26}
Nieskończenie wymiarowa przestrzeń Banacha $\ell^1(\N,\K)$
nie jest refleksywna.
\end{wniosek}
\begin{proof}
Niech $e_n(j)=1$ dla $j=n$ oraz $e_n(j)=0$
dla $j\ne n$. Mamy $\|e_n\|_1=1$ i
$\|e_n-e_m\|_1=2$ dla $n\ne m$, więc ciąg jest
ograniczony, lecz żaden jego podciąg nie jest
ciągiem Cauchy'ego w normie.
Gdyby $\ell^1$ była refleksywna,
Wniosek~\ref{cor:bounded-subsequence26} zapewniałby podciąg
słabo zbieżny do pewnego $x\in\ell^1$.
Twierdzenie~\ref{thm:schur-l1-26} dawałoby jego
zbieżność normową, a więc własność Cauchy'ego.
Sprzeczność dowodzi nierefleksywności $\ell^1$.
\end{proof}

\begin{przyklad}[Słabo zbieżny ciąg bez zbieżności normowej]
\label{ex:en-hilbert26}
Niech $H$ będzie nieskończenie wymiarową przestrzenią Hilberta.
Istnieje w niej nieskończony ciąg ortonormalny $(e_n)$:
po wyborze $e_1,\ldots,e_n$ ich liniowa powłoka jest
właściwą domkniętą podprzestrzenią, więc z rozkładu
ortogonalnego (Twierdzenie~\ref{thm:orth-decomposition13})
można wybrać wektor jednostkowy $e_{n+1}$
w jej dopełnieniu ortogonalnym. Nie zakładamy
ośrodkowości $H$ ani zupełności tego układu ortonormalnego.
Dla każdego $x\in H$ nierówność Bessela
(Lemat~\ref{lem:bessel13}) daje
$\sum_{n\ge1}|\langle e_n,x\rangle|^2\le\|x\|^2$,
więc $\langle e_n,x\rangle\to0$.
Z twierdzenia Riesza
(Twierdzenie~\ref{thm:riesz-hilbert13}) każdy $f\in H'$
ma postać $f(y)=\langle y,x\rangle$ dla pewnego $x\in H$.
Zatem $f(e_n)\to0$ dla każdego $f\in H'$, czyli
$e_n\to0$ słabo w $H$. Jednak $\|e_n\|=1$
dla wszystkich $n$.
Dlatego nieskończenie wymiarowa przestrzeń Hilberta
nie ma własności Schura.
\end{przyklad}

\subsection{Jednostajna wypukłość i twierdzenie Milmana--Pettisa}

\noindent\emph{Powiązania tematyczne w Bibliotece matematycznej:} \atlasref{ME-15::thm:uniform-convexity-15}{ME-15: jednostajna wypukłość przestrzeni funkcji całkowalnych}; \atlasref{ME-15::thm:reflexive-lp-15}{ME-15: refleksywność przestrzeni funkcji całkowalnych}.


\begin{definicja}[Ścisła i jednostajna wypukłość]
\label{def:uniform-convex26}
Niech $X$ będzie przestrzenią unormowaną.
Mówimy, że norma jest \emph{ściśle wypukła}, jeżeli
dla różnych $x,y\in X$ o normach $1$ zachodzi
$\|(x+y)/2\|<1$.
Przestrzeń $X$ jest \emph{jednostajnie wypukła}, jeżeli
dla każdego $\varepsilon\in(0,2]$ istnieje
$\delta(\varepsilon)>0$ takie, że dla wszystkich
$x,y\in X$ o normach $1$ i $\|x-y\|\ge\varepsilon$
zachodzi $\|(x+y)/2\|\le1-\delta(\varepsilon)$.
Przez $\delta_X(\varepsilon)$ oznaczamy dowolną
dodatnią stałą spełniającą ten warunek, nie większą
niż $1$. Jednostajna wypukłość implikuje ścisłą
wypukłość: dla różnych wektorów jednostkowych
wystarczy przyjąć $\varepsilon=\|x-y\|>0$.
\end{definicja}

\begin{przyklad}[Moduł wypukłości przestrzeni Hilberta]
\label{ex:modulus-hilbert26}
Niech $H$ będzie rzeczywistą lub zespoloną przestrzenią
Hilberta. Dla $x,y\in H$ o normach $1$ tożsamość
równoległoboku (Stwierdzenie~\ref{prop:parallelogram13}) daje
$\|x+y\|^2+\|x-y\|^2=4$, czyli
$\|(x+y)/2\|^2=1-\|x-y\|^2/4$.
Jeżeli $\|x-y\|\ge\varepsilon$, to
$\|(x+y)/2\|\le\sqrt{1-\varepsilon^2/4}$.
Zatem $H$ jest jednostajnie wypukła, a można przyjąć
$\delta_H(\varepsilon)=1-\sqrt{1-\varepsilon^2/4}>0$
dla $0<\varepsilon\le2$.
Jeżeli wymiar rzeczywisty $H$ wynosi co najmniej $2$,
jest to optymalna stała: dla dwóch wektorów
ortonormalnych względem rzeczywistego iloczynu skalarnego
$u,v$ bierzemy
$x=\sqrt{1-\varepsilon^2/4}\,u+(\varepsilon/2)v$
i $y=\sqrt{1-\varepsilon^2/4}\,u-(\varepsilon/2)v$.
Oba są jednostkowe i $\|x-y\|=\varepsilon$.
W przypadku zespolonym używamy rzeczywistej części
iloczynu skalarnego, która wyznacza tę samą normę.
Natomiast przestrzeń $\ell^1$ nie jest ściśle wypukła:
$\|e_1\|_1=\|e_2\|_1=\|(e_1+e_2)/2\|_1=1$.
\end{przyklad}

\begin{twierdzenie}[Milman--Pettis]
\label{thm:milman-pettis26}
Każda jednostajnie wypukła przestrzeń Banacha jest refleksywna.
\end{twierdzenie}
\begin{proof}
Wystarczy wykazać, że każdy $\Phi\in X''$ o normie $1$
należy do $J_X(X)$: element zerowy jest obrazem zera,
a dla $\Phi\ne0$ stosujemy wynik do $\Phi/\|\Phi\|$
i mnożymy przeciwobraz przez $\|\Phi\|$.
Obejmuje to również przypadek $X=\{0\}$.
Ustalmy więc $\Phi$ o normie $1$.
Z Twierdzenia~\ref{thm:goldstine26} i charakteryzacji
domknięcia za pomocą sieci istnieje sieć
$(x_\alpha)\subseteq B_X$ taka, że
$J_Xx_\alpha\to\Phi$ słabo-*.
Ustalmy $\varepsilon\in(0,2]$ i wybierzmy
$\delta>0$ mniejsze od $\min\{\delta_X(\varepsilon/2),\varepsilon/4\}$.
Z definicji normy $\Phi$ wybieramy $g\in B_{X'}$
tak, że $|\Phi(g)|>1-\delta/2>0$.
Dzieląc przez $\|g\|$ i mnożąc przez skalar modułu $1$,
otrzymujemy $f\in X'$ o normie $1$ spełniający
$\operatorname{Re}\Phi(f)>1-\delta/2$.
Zbieżność słaba-* daje indeks $\alpha_0$ taki, że
$|f(x_\alpha)-\Phi(f)|<\delta/2$ dla wszystkich
$\alpha\ge\alpha_0$. Wówczas
$\operatorname{Re}f(x_\alpha)>1-\delta>0$.
Ponieważ $\|f\|=1$ i $x_\alpha\in B_X$, mamy
$1-\delta<\|x_\alpha\|\le1$ i dla tych indeksów
możemy położyć $u_\alpha=x_\alpha/\|x_\alpha\|$.
Wektory $u_\alpha$ mają normę $1$, a dzielenie
dodatniej części rzeczywistej przez liczbę z $(0,1]$
daje $\operatorname{Re}f(u_\alpha)>1-\delta$.
Dla $\alpha,\beta\ge\alpha_0$ otrzymujemy
\[
 \left\|\frac{u_\alpha+u_\beta}{2}\right\|
 \ge\operatorname{Re}f\left(\frac{u_\alpha+u_\beta}{2}\right)
 >1-\delta.
\]
Dzięki wybranemu $\delta$ jednostajna wypukłość daje
$\|u_\alpha-u_\beta\|<\varepsilon/2$.
Ponadto $\|u_\alpha-x_\alpha\|=1-\|x_\alpha\|<\delta$,
zatem $\|x_\alpha-x_\beta\|<\varepsilon/2+2\delta<\varepsilon$.
Stąd $(x_\alpha)$ jest siecią Cauchy'ego w normie.
Zupełność $X$ daje $x_\alpha\to x$ w normie:
w przestrzeni metrycznej zupełność dla ciągów
implikuje zupełność dla sieci Cauchy'ego.
Istotnie, dla każdego $n$ wybieramy indeks,
od którego odległości dwóch wyrazów są mniejsze
niż $2^{-n}$; korzystając z kierunkowości, możemy
wybrać rosnące indeksy $\alpha_n$ o tej własności.
Ciąg $(x_{\alpha_n})$ jest Cauchy'ego i zbiega
do pewnego $x$. Dla $\alpha\ge\alpha_n$
mamy $\|x_\alpha-x_{\alpha_n}\|<2^{-n}$,
więc cała sieć zbiega do $x$.
Izometryczność $J_X$ daje zbieżność normową,
a zatem słabą-* sieci $J_Xx_\alpha$ do $J_Xx$.
Z jednoznaczności granicy w topologii słabej-*
Hausdorffa wynika $\Phi=J_Xx$.
Wraz ze skalowaniem z początku dowodu daje to
surjektywność $J_X$.
\end{proof}

\begin{twierdzenie}[Własność Radona--Riesza]
\label{thm:radon-riesz26}
Niech $X$ będzie jednostajnie wypukłą przestrzenią unormowaną,
a $(x_n)\subseteq X$ i $x\in X$. Jeżeli $x_n\to x$
słabo oraz $\|x_n\|\to\|x\|$, to $\|x_n-x\|\to0$.
\end{twierdzenie}
\begin{proof}
Gdy $x=0$, teza wynika wprost ze zbieżności norm.
W przeciwnym razie, ponieważ $\|x_n\|\to\|x\|>0$,
dla dostatecznie dużych $n$ możemy położyć
$u_n=x_n/\|x_n\|$ i $u=x/\|x\|$.
Mamy $\|u_n\|=\|u\|=1$ i dla każdego $f\in X'$
zachodzi $f(u_n)=f(x_n)/\|x_n\|\to f(x)/\|x\|=f(u)$.
Zatem $u_n\to u$ słabo.
Gdyby $\|u_n-u\|\not\to0$, istniałyby
$\varepsilon\in(0,2]$ i podciąg z $\|u_{n_k}-u\|\ge\varepsilon$.
Z jednostajnej wypukłości wynika
$\|(u_{n_k}+u)/2\|\le1-\delta_X(\varepsilon)$.
Tymczasem $(u_{n_k}+u)/2\to u$ słabo, a norma jest
słabo półciągła z dołu, ponieważ
$\|z\|=\sup_{f\in B_{X'}}|f(z)|$
(Wniosek~\ref{cor:norming}) jest supremum
funkcji słabo ciągłych. Otrzymalibyśmy
\[
 1=\|u\|\le\liminf_k\|(u_{n_k}+u)/2\|
 \le1-\delta_X(\varepsilon),
\]
sprzeczność. Stąd $\|u_n-u\|\to0$.
Ostatecznie
$\|x_n-x\|\le\|x_n\|\|u_n-u\|
+|\|x_n\|-\|x\||\to0$, gdyż ciąg norm jest ograniczony.
Zupełność $X$ nie jest potrzebna w tym dowodzie.
\end{proof}

\subsection{Granice i kontrprzykłady}

\begin{twierdzenie}[Słabe domknięcie sfery jednostkowej]
\label{thm:sphere-weak-dense26}
Niech $X$ będzie nieskończenie wymiarową przestrzenią
unormowaną. Wówczas słabe domknięcie sfery jednostkowej
$S_X=\{x\in X:\|x\|=1\}$ jest równe $B_X$.
\end{twierdzenie}
\begin{proof}
Jest to Twierdzenie~\ref{thm:sfera-slaba}; przypomnijmy
argument w kontekście geometrii kuli.
Kula $B_X$ jest wypukła i normowo domknięta,
a więc słabo domknięta
(Wniosek~\ref{cor:weak-convex-closure}).
Ponieważ $S_X\subseteq B_X$, mamy
$\overline{S_X}^{\,w}\subseteq B_X$.
Punkty sfery należą do jej domknięcia.
Niech więc $x\in B_X$ i $\|x\|<1$.
W dowolnym słabym otoczeniu $x$ istnieje
otoczenie podstawowe
\[
 U=\{y\in X:|f_j(y-x)|<\eta,\ j=1,\ldots,m\},
\]
dla pewnych $f_1,\ldots,f_m\in X'$ i $\eta>0$.
Liniowe odwzorowanie $y\mapsto(f_1(y),\ldots,f_m(y))$
do $\K^m$ nie może być różnowartościowe, gdyż
$X$ jest nieskończenie wymiarowa. Istnieje zatem
niezerowy $v\in\bigcap_{j=1}^m\ker f_j$.
Dla każdego $t\ge0$ mamy $x+tv\in U$, gdyż
$f_j(x+tv-x)=0$.
Funkcja $t\mapsto\|x+tv\|$ jest ciągła
(z odwrotnej nierówności trójkąta), w zerze
przyjmuje wartość $\|x\|<1$, a ponadto
$\|x+tv\|\ge t\|v\|-\|x\|\to+\infty$.
Z twierdzenia o wartości pośredniej otrzymujemy
$t_0>0$ z $\|x+t_0v\|=1$.
Zatem $x+t_0v\in U\cap S_X$.
Każde słabe otoczenie $x$ przecina sferę,
co daje $B_X\subseteq\overline{S_X}^{\,w}$
i kończy dowód.
\end{proof}

\begin{uwaga}[Co wynika z twierdzenia Eberleina--\v Smuliana]
\label{rem:weak-compact-what26}
Twierdzenie~\ref{thm:eberlein-smulian26} charakteryzuje
względną słabą zwartość: $\overline A^{\,w}$ jest
słabo zwarte wtedy i tylko wtedy, gdy każdy ciąg
w $A$ ma podciąg słabo zbieżny do punktu $X$.
Nie wymaga się tutaj, aby granica należała do $A$.
Na przykład $A=(0,1)\subseteq\R$ spełnia ten warunek,
ale nie jest zwarty: ciąg $1/(n+1)$ nie ma podciągu
z granicą w $A$.

Dla samego zbioru $A$ słaba zwartość jest równoważna
temu, że każdy ciąg w $A$ ma podciąg słabo zbieżny
do punktu $A$. Jeżeli $A$ jest słabo zwarty, jest
słabo domknięty (topologia słaba jest Hausdorffa);
twierdzenie Eberleina--\v Smuliana daje wymagany
podciąg z granicą w $\overline A^{\,w}=A$.
Odwrotnie, założenie o podciągach z granicą w $A$
zapewnia względną słabą zwartość. Dla dowolnego
$x\in\overline A^{\,w}$ zbiór $A$ jest zawarty
w słabo zwartej przestrzeni $\overline A^{\,w}$.
Samo twierdzenie o względnej zwartości nie pozwala
jednak zastąpić domknięcia domknięciem ciągowym
w dowolnej przestrzeni topologicznej.
W naszym przypadku stosujemy konstrukcję z kroku
trzeciego dowodu Twierdzenia~\ref{thm:eberlein-smulian26}
do $\Phi=J_Xx$: skonstruowany ciąg w $A$ ma
podciąg słabo zbieżny do $y\in A$. Punkt $y$ jest
słabym punktem skupienia tego ciągu, więc ten sam
argument normujący daje $\Phi=J_Xy$.
Z różnowartościowości $J_X$ mamy $x=y\in A$.
Stąd $A=\overline A^{\,w}$ i $A$ jest słabo zwarty.

Słaba zwartość nie oznacza zwartości normowej.
W nieskończenie wymiarowej przestrzeni Hilberta
kula $B_H$ jest słabo zwarta na mocy
Twierdzeń~\ref{thm:milman-pettis26}
i~\ref{thm:reflexive-ball26}, ale zawiera ciąg
ortonormalny z Przykładu~\ref{ex:en-hilbert26}.
Dla $n\ne m$ mamy $\|e_n-e_m\|^2=2$, więc żaden
jego podciąg nie jest normowo zbieżny.
Z kolei w $\ell^1$ sfera ma zero w słabym
domknięciu na mocy Twierdzenia~\ref{thm:sphere-weak-dense26},
ale żaden ciąg jej punktów nie zbiega słabo do zera,
gdyż własność Schura wymuszałaby zbieżność norm do zera.
W przestrzeniach innych niż przestrzenie Banacha
równoważności zwartości słabej i słabej zwartości
ciągowej nie wolno zakładać bez dodatkowego dowodu.
\end{uwaga}

\begin{przyklad}[Refleksywność bez jednostajnej wypukłości danej normy]
\label{ex:reflexive-not-uniform26}
Niech $X=\K^2$ z normą $\|(a,b)\|_\infty=\max\{|a|,|b|\}$.
Jest to przestrzeń Banacha, a jej kula jednostkowa
jest normowo zwarta jako domknięty i ograniczony
zbiór w skończonym wymiarze. Jest więc również
słabo zwarta, ponieważ identyczność z topologii
normowej do słabej jest ciągła.
Twierdzenie~\ref{thm:reflexive-ball26} daje refleksywność.
Jednak dla $u=(1,1)$ i $v=(1,-1)$ zachodzi
$\|u\|_\infty=\|v\|_\infty=1$,
$\|u-v\|_\infty=2$ oraz $\|(u+v)/2\|_\infty=1$.
Norma nie jest nawet ściśle wypukła, a więc nie
jest jednostajnie wypukła.
Przykład dotyczy ustalonej normy: ta sama przestrzeń
z normą euklidesową jest jednostajnie wypukła.
\end{przyklad}
